\documentclass[11pt]{article}
\usepackage[T1]{fontenc}
\usepackage[utf8]{inputenc}
\usepackage{amsmath}
\usepackage{booktabs}
\usepackage{graphicx}
\usepackage{hyperref}

\newcommand{\term}[1]{\textit{#1}}
\newcommand{\ctu}{content translation unit}
\newcommand{\toolname}{localization engine}

\title{Keeping Structure Intact When Localizing LaTeX}
\author{Technical Notes}
\date{}

\begin{document}
\maketitle

\begin{abstract}
Teams that maintain bilingual technical documents in \LaTeX{} face a recurring
trade-off: machine translation is fast, but naive string replacement destroys
equations, cross-references, and floats. This note argues for treating
structure as a first-class constraint during localization. We sketch a
simple integrity model, summarize practical failure modes, and outline a
workflow in which prose is translated while markup remains under
deterministic control.
\end{abstract}

\section{Motivation}
\label{sec:motivation}

A compiled paper is more than running text. Authors rely on stable labels,
numbered equations, and captions that continue to point at the right
objects after revision. When those documents must appear in a second
language, the temptation is to paste whole files into a chat model and hope
for the best. The model may produce fluent prose and still break the
document: a missing brace, a rewritten \texttt{\textbackslash ref}, or a
caption that no longer matches its float.

Maintainers therefore need a different contract. The unit of work should be
a \term{\ctu{}}---typically a paragraph or caption---while macros, math, and
reference tokens are masked before the model sees the text and restored
afterward.\footnote{Human review remains essential for terminology and tone;
integrity checks only guarantee that structure survived the machine pass.}
In practice this means a local \term{\toolname{}} can call an
OpenAI-compatible endpoint without surrendering compile safety.

Typical pressure points include:

\begin{itemize}
  \item Equations that must stay byte-identical across languages.
  \item Captions whose wording changes while float labels must not.
  \item Citations and page references that resolve only if keys survive.
\end{itemize}

Public discussion of continuous localization often focuses on strings in
software UIs; the same integrity concerns apply to scholarly \LaTeX{}, as
noted in recent operator notes~\cite{integrity2024}. For background on
incremental memory-backed workflows, see~\cite{workflow2025}. Online
tooling overviews are easy to find; one starting point is the
\href{https://www.latex-project.org/}{\LaTeX{} Project} site.

\section{An integrity sketch}
\label{sec:model}

Consider a source fragment that mixes prose with inline mathematics such as
$E = mc^{2}$ and $\alpha + \beta = \gamma$. A safe pipeline extracts the
linguistic span, replaces protected tokens with opaque placeholders, and
asks the model only for the span. After validation, placeholders are mapped
back so that display material is unchanged:

\begin{equation}
\label{eq:energy}
E = mc^{2}
\end{equation}

Equation~\eqref{eq:energy} is deliberately elementary. The point is not the
physics; it is that the rebuilt file must still contain the same tokens the
compiler expects. Neighboring paragraphs may supply context to the model,
but they must not become an invitation to rewrite labels or environments.

\section{Observability for maintainers}
\label{sec:observe}

Operators need coarse signals before they accept a batch. Table~\ref{tab:signals}
lists a minimal set of checks that catch the most expensive failures early.

\begin{table}[ht]
  \centering
  \begin{tabular}{@{}ll@{}}
    \toprule
    Signal & Why it matters \\
    \midrule
    Placeholder multiset & Detects dropped or invented structure tags \\
    Status gates & Prevents shipping unfinished machine drafts \\
    Human locks & Protects reviewed wording from auto-overwrite \\
    \bottomrule
  \end{tabular}
  \caption{Lightweight checks that keep localization batches compile-safe.}
  \label{tab:signals}
\end{table}

The same discipline applies to figures. Figure~\ref{fig:pipeline} stands in
for a pipeline diagram: prose in captions is fair game for translation;
layout boxes and labels are not.

\begin{figure}[ht]
  \centering
  \rule{0.72\linewidth}{2.4cm}
  \caption{Schematic of a structure-preserving pass: parse, mask, translate
    prose, validate, and rebuild. The shaded band represents protected
    tokens that never enter the model as editable text.}
  \label{fig:pipeline}
\end{figure}

After a rebuild, a maintainer should still be able to jump from
Section~\ref{sec:motivation} to Equation~\eqref{eq:energy} and to the summary
in Table~\ref{tab:signals} on page~\pageref{tab:signals}.

\section{Working vocabulary}
\label{sec:vocab}

A short shared vocabulary reduces review friction:

\begin{description}
  \item[Segment] An atomic \ctu{} stored in translation memory with a stable
    identifier and status.
  \item[Placeholder] A masked token that stands for math, a reference, or
    another protected construct during machine translation.
\end{description}

Together with the checks in Table~\ref{tab:signals}, this vocabulary keeps
conversations focused on meaning rather than on recovering broken markup.

\section{Closing}
\label{sec:close}

Localizing \LaTeX{} is feasible when integrity is enforced before fluency is
celebrated. Masking, validation, and human priority turn a risky paste into
a repeatable workflow. Teams that adopt that order spend less time repairing
documents and more time improving the prose itself---which is the only part
the reader was meant to notice.

\begin{thebibliography}{9}
\bibitem{integrity2024}
  A.~Maintainers.
  \newblock Integrity over fluency in structured localization.
  \newblock \emph{Technical Note}, 2024.

\bibitem{workflow2025}
  A.~Maintainers.
  \newblock Sync, translate, and rebuild without rewriting macros.
  \newblock \emph{Operator Notes}, 2025.
\end{thebibliography}

\end{document}
